\section{Unique Games with low completeness}\label{sec:johnson_2}
In this section we give an analysis which works for UG instances with arbitrary small completeness (but bounded away from $0$). The only step that changes is concluding that there is a subcube with large shift-potential (after conditioning). 

Recall that in the proof of the analogous lemma for large completeness (Lemma~\ref{lem:sp-j}) 
we used Claims~\ref{claim:subcube-expanse} and \ref{claim:potentials} to conclude that if a subcube has large global shift-partition potential it also has large shift-partition potential. This was possible 
because all the vertices with high global value will also have high value inside the subcube as the expansion of the subcube is small ($O(\eps)$ if the completeness was $1-\eps$). This fact is no longer true in the low completeness regime, as subcubes 
now may have expansion close to $1$.
To circumvent this issue, we no longer go via the intermediate global shift-partition route and 
directly try to conclude that the shift-partition potential is large on a subcube. Towards this end 
we need a stronger conclusion of global hypercontractivity which we refer to 
as an ``edge-covering'' statement: the subcubes on which some part of the shift partition is large cover nearly all of the the internal edges of the shift-partition. 

As we shall see, if the value of our pseudodistribution is originally $c$, the 
fraction of internal edges is at least $c^2$ (in particular $c^2$-fraction of edges are satisfied by both $X$ and $X'$), and as we are able to cover almost all internal edges we get that there exists a subcube with large completeness and in fact large shift-partition potential. Given this version of Lemma~\ref{lem:sp-j}, the rest of the analysis remains the same. 

\paragraph{Notation:}
For the purposes of this section we first consider the simpler shift-partition defined by functions $\{G_s: V(G) \rightarrow \R[X,X']\}_{s \in \Sigma}$:
\[G_s(u) = \one(X_u - X'_u = s).\]
Note that $G_s$'s cover all the vertices, that is $\sum_s \delta(G_s) = 1$ and satisfy Booleanity: $G_s(u)^2 = G_s(u)$ for all $u \in V(G)$.

Given two assignments $X,X'$, let $\val_u(X \wedge X')$ denote the fraction of edges incident on $u$ that are satisfied by \emph{both} $X$ and $X'$. Note that $\val_u(X \wedge X')$ is a degree four polynomial in $(X,X')$: 
\[\val_u(X \wedge X') = \E_{v: (u,v) \in E}[(\sum_{\sigma \in \Sigma} X_{u,\sigma}X_{v,\pi_{u,v}(\sigma)})(\sum_{\sigma \in \Sigma} X'_{u,\sigma}X'_{v,\pi_{u,v}(\sigma)})].\]
Let $\val_I(X \wedge X') = \E_u[\val_u(X \wedge X')]$ denote the fraction of edges in the graph that are satisfied by both $X$ and $X'$. Further for an $i$-restriction $a \in {[n] \choose i}$ and a vertex $u \in J|_a$ let $\val^a_u(X \wedge X')$ denote the fraction of edges incident on $u$ in $J|_a$ that are satisfied by both $X$ and $X'$. This is also a degree $4$ polynomial in $(X,X')$.

\begin{definition}[Approximate Indicator for Dense Subcubes]
Let $\eps_0,\ldots,\eps_r \in (0,1)$ to be determined later. For any $i$-restriction $a$ let $R_{s,a}$ denote the (approximate) indicator that $G_s$ is $\eps_i$-dense in $J|_a$, but not $\eps_{|b|}$-dense inside any subcube $J|_b$ where $b \subsetneq a$. To define this as a polynomial let us first define the approximate indicator $p_{<\beta,\nu}(x)$ as $1 - p_{\beta,\nu}(x)$. One can check that $p_{<\beta,\nu}$  approximates $\Ind[x < \beta]$ when $x \in [0,1]$ with similar properties as $p_{\beta,\nu}$ (Theorem~\ref{thm:step-approx}) that approximates $\Ind[x \geq \beta]$ for $x \in [0,1]$. Formally using the polynomial approximation for an indicator define 
\[R_{s,a} = p_{\eps_i,\nu}(\delta(G_s|_a))\prod_{j < i}\prod_{b \subset a: |b| = j}p_{<\eps_j,\nu}(\delta(G_s|_b)).\] 
\end{definition}


\subsection{The Edge Covering Theorem}
Our argument will need an upper bound on the number of
events $R_{s,a}$ that can occur simultaneously, which 
roughly speaking asks how many $i$-restricted subcubes can a given set $F$ be dense on. 
As stated, there is no good upper bound for this: 
if $F$ is dense on an $i-1$-restricted subcube
then it would be quite dense on many $i$-restricted subcubes containing it. This is the reason 
that in the event $R_{s,a}$, we required that the $i$-restricted subcube is dense but there is no $j$-restricted subcube ($j < i$) containing it on which our set is still somewhat dense.

\begin{claim}\label{claim:dense-subcubes}
If $\eps_0, \ldots, \eps_r$ satisfy that $\eps_{i-1} \leq \eps_i/(2^{i+1} i)$ then we get that:
\[\vdash_{\ell^{O(\ell^i/\nu)}} \E_{a \sim {[n] \choose i}}[R_{s,a}] \leq \frac{4\delta(G_s)}{\eps_i^2\ell^i} + O(\nu).\]
\end{claim}
\begin{proof}
We first present an argument when the polynomials $p_{\beta,\nu}$ in $R_{s,a}$ are replaced by indicators. We then give a sketch of how to convert the proof into an SoS proof. 

Converting the $p_{\beta,\nu}$'s to indicators, we get that $R_{s,a}$ corresponds to $T_{s,a}$: 
\begin{align*}
    T_{s,a} 
    &= \one(\delta(G_s|_a) \geq \eps_i) \prod_{j < i}\prod_{b \subset a: |b| 
    = j}\one(\delta(G_s|_b) < \eps_j)
\end{align*}
that is the event that $G_s$ is $\eps_i$-dense in $a$, but for all subsets $b$ of $a$, it is at most $\eps_j$-dense. We will show that the fraction of such restrictions $a$ must be small.

We consider the Fourier-analytic function $f_{i,F}$ defined in Section~\ref{sec:fourier}, Definition~\ref{def:f_i} and use its alternative formula from Lemma~\ref{lem:restriction}:
\[f_{i,F}(a) = \sum_{b \subseteq a} (-1)^{i -|b|}\delta(F|_b) \geq \delta(F|_a) - \sum_{b \subsetneq a} \delta(F|_b).\]
Technically we have only defined these functions for the Cayley-version of the Johnson graph $C_{n,\ell,\alpha}$, but one can use the above definition for Johnson graphs and derive the same properties that we use here upto $o(1)$ error terms. We will ignore these $o(1)$-error terms in this proof.

If $T_{s,a} = 1$ we get that, $\delta(F|_a) \geq \eps_i$, but each $\delta(F|_b)$ is at most $\eps_{|b|} < \eps_i/(2^{i+1} i)$ so we get:
\[f_{i,G_s}(a) \geq \eps_i - \sum_{j < i} {i \choose j}\eps_j \geq \eps_i - \sum_{j < i} 2^i \eps_j \geq \frac{\eps_i}{2}.\]
This immediately gives: 
\begin{equation}\label{eq:rsa1}
T_{s,a} \leq 1 \leq 4f_{i,G_s}(a)^2/\eps_i^2,
\end{equation}
if $T_{s,a} = 1$. If $T_{s,a} = 0$ then the above continues to hold.

On the other hand by the definition of $f_{i}$ Lemma~\ref{lem:3.2} gives that, 
\[\E[f_{i,G_s}(a)^2] = \frac{W^i[G_s]}{\ell^i} \leq \frac{\E[G_s^2]}{\ell^i}.\]
Combining with equation~\ref{eq:rsa1} we get:
\[\E_{a \sim {[n] \choose i}}[T_{s,a}] \leq \E_{a}\left[\frac{4f_{i,G_s}(a)^2}{\eps_i^2}\right] \leq \frac{4\delta(G_s)}{\eps_i^2\ell^i}.\]

\paragraph{SoS-ing the proof:} We sketch an SoS proof for the statement for $R_{s,a}$ with $a \in {[n] \choose 1}$. The full statement follows analogously. 

Let $\delta_0$ denote $\delta(G_s)$ and $\delta_1$ denote $\delta(G_s|_a).$ Analogous to \eqref{eq:rsa1} we will first show that:
\begin{equation}\label{eq:rsa-non-sos}
\frac{4(\delta_1 - \delta_0)^2}{\eps_1^2} + 2\nu - R_{s,a} > \nu
\end{equation}
if $\delta_1,\delta_0 \in [0,1]$. This will not be an SoS proof, but after showing this we can use a blackbox theorem to convert it into an SoS proof by the properties of our domain $[0,1]^2$. \eqref{eq:rsa-non-sos} follows by case analysis. 
\begin{enumerate}
    \item If $R_{s,a} > \nu$ it implies that $p_{\eps_1,\nu}(\delta_1) > \nu$ and $p_{<\eps_0,\nu}(\delta_0) > \nu$. By the definition of the polynomials $p_{\beta,\nu}$ this implies that $\delta_1 > \eps_1$ and $\delta_0 < \eps_0 < \eps_1/2$, which gives that $(\delta_1 - \delta_0)^2 \geq \eps_1^2/4$. Since $R_{s,a}$ is always $\leq 1$ we get:
    \[R_{s,a} \leq 1 \leq \frac{4(\delta_1 - \delta_0)^2}{\eps_1^2} < \frac{4(\delta_1 - \delta_0)^2}{\eps_1^2}+\nu.\]
    \item If $R_{s,a} \leq \nu$ we get:
    \[R_{s,a} \leq \frac{4(\delta_1 - \delta_0)^2}{\eps_1^2} + \nu.\]
\end{enumerate}
Hence rearranging the above and adding $\nu$ on both sides gives \eqref{eq:rsa-non-sos}. We can now apply Theorem~\ref{thm:blackbox-sos} that shows that positive polynomials bounded away from $0$ on the domain $[0,1]^k$ have a bounded degree SoS proof to get:
\[\delta_0,\delta_1 \in [0,1] \vdash_{\exp(1/\nu)} R_{s,a} \leq \frac{4(\delta_1 - \delta_0)^2}{\eps_1^2} + 2\nu.\]

One can check that the rest of the proof is already SoS therefore we get:
\[\cA_I \vdash_{\tO(1/\nu)} \delta_0,\delta_1 \in [0,1] \vdash_{\exp(1/\nu)} \E_{a \sim {[n] \choose 1}}[R_{s,a}] \leq \frac{4\delta(G_s)}{\eps_1^2\ell} + 2\nu.\]

One can do the exact same steps and use the SoS degree bound in Theorem~\ref{thm:blackbox-sos} to get an SoS proof of the full statement with degree $\ell^{O(\ell^i/\nu)}$. The bound that comes from Theorem~\ref{thm:blackbox-sos} is exponential in the degree of the polynomial $f$ that we want to prove is SoS, and hence our degree bound is also $\exp(\ell^i/\nu)$. We believe there could be a more direct analysis to get a degree bound of $\poly(1/\nu, \ell^i)$ instead.
\end{proof}

We now move on to stating the edge covering theorem. 
Below, we inspect the edges that stay inside a part in the shift partition, and show that they can nearly be encapsulated within subcubes on which some part becomes dense.
\begin{theorem}[Edge-Covering Theorem]\label{thm:edge-covering-restated}
For all constants $r \in \N$, $\eps_0, \ldots, \eps_r \in [0,1]$, where $\eps_r \leq \exp(-r)$ and $\eps_0, \ldots, \eps_r$ satisfy that $\eps_{i-1} \leq \eps_i^5/(2^{6r})$, we get that:
\begin{equation}\label{eq:edge-cov-final}
\cA_I \vdash_{\ell^{\poly(\ell^r/\nu)}}  \val_I(X \wedge X') \leq T_0(X,X') + \ldots + T_r(X,X') + \sf{err},
\end{equation}
where:
\begin{align*}
&T_0(X,X') = \sum_s p_{\eps_0,\nu}(\delta(G_s))\E_{u \sim G}[ G_s(u)\val_u(X \wedge X')],\\
&T_i(X,X') = \ell^{i}\E_{a \in {[n] \choose i}}[\sum_s R_{s,a}
\E_{u \sim J|_a}[G_s(u)\val^a_u(X \wedge X')]] \\
&{\sf err} = 4(1-\alpha)^{r+1} + 2^{6r}\max_{i}\frac{\eps_{i-1}}{\eps_i^4}+ |\Sigma|\ell^{O(r)}\sqrt{\nu}.
\end{align*}
\end{theorem}

\begin{proof}
We will first give a proof of this statement with the definition of $R_{s,a}$ replaced by actual indicator functions in place of the polynomials $p_{\beta,\nu}$. Formally, for an $i \leq r$-restriction $a \in {[n] \choose i}$ let $R_{s,a}$ be an indicator variable that is 1 if $\delta(G_s|_a) \geq \eps_i$ and $\delta(G_s|_b) \leq \eps_{|b|}$ for all $b \subsetneq a$. We will then give a sketch of how to convert this into an \sos proof and work with the true definition of $R_{s,a}$.

For any set $S$ let $\#sat(S)$ denote the number of edges inside $S$ that are satisfied by both $X$ and $X'$. We first have that:
\begin{equation}\label{eq:edge_cover0}
\val_{I}(X \wedge X') \leq \sum_s \frac{\#sat(G_s)}{|E|},
\end{equation}
since we know (from Fact~\ref{fact:z-vars}) that any edge that crosses between parts in the shift-partition must be violated by at least one of $X$ or $X'$. Let us therefore upper bound $\#sat(G_s)$.


Consider the set indicated by $G_s$ for fixed $s \in \Sigma$. Let the set of dense subcubes for $G_s$ be defined as follows:
\begin{align*}
&\cC^0_s = \begin{cases}\{J\} ~~~~~~~~~\text{if $\delta(G_s) \geq \eps_0$} \\
\phi ~~~~~~~~~~~~\text{otherwise}
\end{cases} \\
&\cC^{(1)}_s = \left\{a \in [n]: R_{s,a} = 1\right\} \\
&\vdots \\
&\cC^{(r)}_s = \left\{a \in {[n] \choose r}: R_{s,a} = 1\right\}. 
\end{align*}


Let $D_s$ denote the dense part of $G_s$: $\cup_{i \leq r} \cup_{a \in \cC^{(i)}_s} G_s \cap J|_a$ and let $H_s = G_s \setminus D_s$ be the pseudorandom part of $G_s$. It is easy to check that $H_s$ is $(r, \eps_r)$-pseudorandom. The edges inside $G_s$ can be divided into four types of edges: edges that stay inside the dense subcubes, edges that go between two different dense subcubes, edges inside $H_s$ and edges that go between $H_s$ and $D_s$. Applying this subdivision of edges on $\#sat(G_s)$ we get that:
\begin{align}
\frac{\#sat(G_s)}{|E(J)|} 
&\leq \sum_{i \leq r} \sum_{a \in \cC^{(i)}_s} \frac{\#sat(G_s \cap J|_a)}{|E(J)|}  +
\frac{|E(H_s, H_s)|}{|E(J)|} + \frac{|E(H_s, D_s)|}{|E(J)|}\notag\\\label{eq:edge-cov1}
+ 
&\sum_{i, j \leq r} 
\sum\limits_{h=0}^{\min(i,j)}
\sum_{a \in \cC^{(i)}_s}
\sum_{H\subseteq a, |H| = h}
\sum_{\substack{b \in \cC^{(j)}_s \\ a \neq b, a\cap b = H}} \frac{|E(J|_a, J|_b)|}{|E(J)|}.
\end{align}

We will check that the $i^{th}$-summand in the first term can be rearranged to give the $s^{th}$-summand in $T_i(X,X')$. We have that:
\begin{equation}\label{eq:edge-cov2}
\sum_{a \in \cC^{(i)}_s} \frac{\#sat(G_s \cap J|_a)}{|E(J)|} = \frac{{[n] \choose i}|E(J|_a)|}{|E(J)|}\E_{a \sim {[n] \choose i}}\left[R_{s,a}\frac{\#sat(G_s \cap J|_a)}{|E(J|_a)|}\right].
\end{equation}
By a direct computation we get that $\frac{{[n] \choose i}|E(J|_a)|}{|E(J)|} \leq \ell^i$. As for the second term above we get that, 
\begin{equation*}
\frac{\#sat(G_s \cap J|_a)}{|E(J|_a)|} \leq \frac{\sum_{u \in J|_a}\Ind[u \in G_s]\#sat_a(u)}{|V(J|_a)|d_a} = \E_{u \in J|_a}[\Ind[u \in G_s]\val_u^a(X \wedge X')],
\end{equation*}
where $d_a$ denotes the degree of a vertex inside the subgraph $J|_a$ and $\#sat_a(u)$ denotes the number of edges in $J|_a$ that are incident on $u$ and are satisfied by both $X$ and $X'$. So plugging this into \eqref{eq:edge-cov2} and summing up over $i \leq r$ we get that the first term in \eqref{eq:edge-cov1} gives us:
\begin{equation}\label{eq:first-term-ec}
\sum_{i \leq r} \sum_{a \in \cC^{(i)}_s} \frac{\#sat(G_s \cap J|_a)}{|E(J)|} \leq \sum_{i \leq r} \ell^i \E_{a \sim {[n] \choose i}}[R_{s,a}\E_{u \in J|_a}[\Ind[u \in G_s]\val_u^a(X \wedge X')]]. 
\end{equation}


We will bound the other terms in \eqref{eq:edge-cov1} by some small constant. 

\paragraph{Bounding the last term in~\eqref{eq:edge-cov1}.} 
We focus on $i\leq j$, and in the end multiply the 
bound we get by factor $2$. 

If $R_{s,\phi} = 1$, that is, $\delta(G_s) \geq \eps_0$ then we can check that this is $0$ since none of the other subcubes will have $R_{s,a} = 1$. So let us assume that this is not the case. 
Fix $i$ and $H\subseteq i$ of size $h$. For any $a \in {[n] \choose i}, b \in {[n] \choose j}$
intersecting in a set of size $h$ we have that
$\frac{|E(J|_a, J|_b)|}{|E(J)|} \leq \left(\frac{2\ell}{n}\right)^{i+j-h}$.
Indeed, to see this one can think of sampling an edge 
$(x,y)$ in $J$ uniformly at random, and asking what 
is the probability that $a\subseteq x$ and $b\subseteq y$. For that, we want $a\cap b$ to be contained in 
both $x$ and $y$ (which happens with probability at most $(2\ell/n)^h$), and that $a\setminus(a\cap b)\subseteq x$, $b\setminus(a\cap b)\subseteq y$
which happens with probability at most $(2\ell/n)^{i-h}\cdot (2\ell/n)^{j-h}$. 
Thus, 
we get that the last term in~\eqref{eq:edge-cov1}
is at most
\[
\sum_{i\leq j \leq r} 
\sum\limits_{h=0}^{i-1}
\sum_{a \in \cC^{(i)}_s}
\sum_{H\subseteq a, |H| = h}
|\{b~|b\in \cC^{(j)}_s, a\cap b = H, a\neq b\}|\left(\frac{2\ell}{n}\right)^{i+j-h}.
\]
Here we note that $h$ cannot be $i$, since then we 
would be looking at sub-cubes $J|_a\subseteq J|_b$, 
and so we cannot have both $R_{s,a} = 1$ and $R_{s,b}=1$.

Applying Claim~\ref{claim:dense-subcubes} on the 
restriction $G_s|_H$, we get that 
\[
|\{b~|b\in \cC^{(j)}_s, a\cap b = H, a\neq b\}|
\leq \frac{n^{j-h}}{\ell^{j-h}}\frac{\delta(G_s|_H)}{\eps_{j}^2}
\leq \frac{n^{j-h}}{\ell^{j-h}}\frac{\eps_h}{\eps_{j}^2}.
\]
We used the fact that as $R_{s,b} = 1$,  $G_s$ is at most $\eps_h$ dense in $J|_h$. Plugging that above we get that the last term 
on the right hand side of~\eqref{eq:edge-cov1} is 
at most
\[
\sum_{i\leq j \leq r} 
\sum\limits_{h=0}^{i-1}
\sum_{a \in \cC^{(i)}_s}
\sum_{H\subseteq a, |H| = h}
\frac{n^{j-h}}{\ell^{j-h}}\frac{\eps_h}{\eps_{j}^2}
\leq 
\sum_{i, j \leq r} 
2^i \cdot i\cdot \frac{n^{j-h}}{\ell^{j-h}}\frac{\eps_h}{\eps_{j}^2}
\cdot |\cC^{(i)}_s|\left(\frac{2\ell}{n}\right)^{i+j-h},
\]
and as $|\cC^{(i)}_s|\leq \frac{\delta(G_s)n^i}{\eps_i^2\ell^i}$ due to 
Claim~\ref{claim:dense-subcubes}, we get further upper
bound this by
\[
\sum_{i, j \leq r} 
2^i \cdot i\cdot \frac{n^{i+j-h}}{\ell^{i+j-h}}\frac{\eps_h\delta(G_s)}{\eps_{j}^2\eps_i^2}
\left(\frac{2\ell}{n}\right)^{i+j-h}
\leq
2^{3r+1} r^2\delta(G_s) \max_{i}\frac{\eps_{i-1}}{\eps_i^4}
\leq 2^{6r}\max_{i}\frac{\eps_{i-1}}{\eps_i^4}\delta(G_s).
\]
Summing this over $s$ gives that the contribution of the last term in~\eqref{eq:edge-cov1} to~\eqref{eq:edge_cover0} is at most
$2^{6r}\max_{i}\frac{\eps_{i-1}}{\eps_i^4}$.

% and we next upper bound the size of the last set. 
% Since $a$ and $b$ are distinct, we must either 
% have that $i\neq j$ or that $i=j$ and $h<i$. 
% We first address the first case, and say without
% loss of generality that $i<j$. There 


% For a particular $i,j$ the number of terms in the summand are at most the number of dense $i$-subcubes times the number of dense $j$-subcubes which we can bound by Claim~\ref{claim:dense-subcubes}. 
% Using the upper bound in Claim~\ref{claim:dense-subcubes} we get that:
% \begin{align}
% \sum_{i, j \leq r} \sum_{\substack{a \in \cC^{(i)}_s, \\ b \in \cC^{(j)}_s \\ a \neq b}} \frac{|E(J|_a, J|_b)|}{|E(J)|} &\leq \sum_{0 \neq i,j \leq r} \Ind[\delta(G_s) < \eps_0]
% \left(\frac{2\ell}{n}\right)^{i+j}\cdot \left(\frac{4\delta(G_s)n^i}{\eps_i^2\ell^i} \right)\cdot \left(\frac{4n^j\delta(G_s)}{\eps_j^2\ell^j}\right) \notag\\
% &\leq \frac{O(r^2)\delta(G_s)^2 \Ind[\delta(G_s) < \eps_0]}{\eps_1^4} \notag \\
% &\leq \frac{O(r^2)\eps_0\delta(G_s)}{\eps_1^4}.\label{eq:second-term-ec}
% \end{align}


\paragraph{Bounding the second and third term in \eqref{eq:edge-cov1}.} 
Here, we use Theorem~\ref{thm:structure-johnson}. 
The point is that the set $H_s$ is pseudo-random, 
and hence we can get good control over the edges from it to other sets. 

More precisely, for the second term note that $H_s$ is  $(r,\eps_r)$-pseudorandom  as per Definition~\ref{def:p.r.}, 
so by Theorem~\ref{thm:structure-johnson}) with $\gamma =\eps_r$ we get: 
\begin{equation}\label{eq:third-term-ec}
\frac{|E(H_s)|}{|E(J)|} = \delta(H_s) - \ip{H_s, LH_s} \leq 2\delta(H_s)(1-\alpha)^{r+1} \leq 2\delta(G_s)(1-\alpha)^{r+1},
\end{equation}
where the first inequality holds as long as $\eps_r \leq (1-\alpha)^{r+1}\exp(-r)$. Note that the second term in the structure theorem $q_a(H_s)(\delta(H_s|_a) - \eps_r)$ is non-positive since $q_a$ is \sos and $H_s$ is pseudorandom. The third term corresponding to Booleanity, $B(H_s)$ is $0$ since $H_s^3(u) = H_s(u)$.
Summing up over $s$ yields that the contribution of this to the right hand side of~\eqref{eq:edge_cover0} 
is at most $2(1-\alpha)^{r+1}$.

The bound on the third term in~\eqref{eq:edge-cov1} is similar. We have that
\begin{equation}\label{eq:fourth-term-ec}
\frac{|E(H_s,D_s)|}{|E(J)|} = \ip{D_s, J H_s} \leq \sqrt{\ip{D_s,D_s} \cdot \ip{H_s, J^2 H_s}} \leq 2\delta(G_s)(1-\alpha)^{r+1}.
\end{equation}
In the last inequality, we used $\ip{D_S,D_s}\leq \delta(G_s)$ and 
$\ip{H_s, J^2 H_s}\leq 2\delta(G_s)(1-\alpha)^{2(r+1)}$ 
which follows by Theorem~\ref{thm:structure-johnson} 
exactly as above. Summing up over $s$ yields total contribution of at most $2(1-\alpha)^{r+1}$.


\paragraph{SoS-ing the proof:} We give a proof sketch of how to convert this into an SoS proof. To avoid Booleanity error we will apply \eqref{eq:edge-cov1-sos} multiple times with $R_{s,a}$ replaced with $p_{\beta,\nu}$'s and take an average. We know that $R_{s,a} = p_1(\delta_1) \cdot \ldots\cdot p_t(\delta_t)$ where $p_i$'s are approximate indicators $p_{\beta_i,\nu}$ (or $p_{<\beta_i,\nu}$) and $\delta_i$'s are of the form $\delta(G_s|_b)$ for $b \subseteq a$. Then define $R^c_{s,a}$ as $p^c_1(\delta_1) \cdot \ldots\cdot p^c_t(\delta_t)$ where $p^c_i(\delta_i) = p_{\beta_i+c\nu^2,\nu^2}$ (resp. $p_{<\beta_i+c\nu^2,\nu^2}$) for all $c \in \{0,\ldots,\lfloor 1/\nu \rfloor\}$.

Note that $G_s$'s and the functions $\#sat(G_s \cap J|_a)$ are actual indicators as written since we can express $G_s(u) = \Ind[X_s - X'_u = s]$ and $\Ind[e \text{ is satisfied }]$ as a polynomial in $(X,X')$. Define $H^c_s(u) = G_s(u)\prod_{a \subset u: |a| \leq r}(1-R^c_{s,a})$ and $D^c_{s} = G_s(u)(1 - H^c_{s,a})$. Defining dense subcubes with $R^c_{s,a}$ we get the following analogue of \eqref{eq:edge-cov1}:
\begin{align}\label{eq:edge-cov1-sos}
\frac{\#sat(G_s)}{|E(J)|} &\leq \sum_{i \leq r} \sum_{a \in {[n] \choose i}}R^c_{s,a} \frac{\# sat(G_s \cap J|_a)}{|E(J)|} +
\ip{H^c_s, AH^c_s} + \ip{D^c_s, JH^c_s} \notag\\+&\sum_{i, j \leq r} 
\sum\limits_{h=0}^{\min(i,j)}
\sum_{a \in {[n] \choose i}} R^c_{s,a}
\sum_{H\subseteq a, |H| = h}
\sum_{\substack{b \in {[n] \choose j} \\ a \neq b, a\cap b = H}} R^c_{s,b} \frac{|E(J|_a, J|_b)|}{|E(J)|}.
\end{align}
We then take an expectation over $c \in \{0,\ldots,\lfloor 1/\nu \rfloor\}$ and then a sum over $s \in \Sigma$.

Now we can use the same analysis to upper bound each of these terms. The analysis for the first term follows as is; we get an $\E_c[R^c_{s,a}]$ (instead of $R_{s,a}$) but we will ignore this slight difference as it does not affect our final analysis of the algorithm.
For the second term we use the \sos version of Claim~\ref{claim:dense-subcubes} (as stated) to say that $\E_{a \in {[n] \choose i}}[R^c_{s,a}]$ is small and the same tricks like those in the SoS-ization of Claim~\ref{claim:dense-subcubes}. For the third and fourth terms we apply the structure theorem on the almost-Boolean function $H_s$ and bound the Booleanity term $\E_c[\sum_s B(H^c_s)]$ by $|\Sigma|\ell^{O(r)}\sqrt{\nu}$ instead. Additionally in the fourth term we apply an \sos version of Cauchy-Schwarz. We omit these details as we've carried out similar arguments before, but formally prove the upper bound on $\E_c[\sum_s B(H^c_s)]$ in Claim~\ref{claim:bool} in the appendix.
\end{proof}


\subsection{The Main Lemma}
In this section, we consider UG instances that have completeness bounded away from $1$ and possibly very close to $0$. The analysis of our algorithm differs only in one step: concluding that the shift potential is large. 


The following claim asserts that taking two assignments $X, X'$ from our pseudo-distribution, 
the expected number of edges that are satisfied by 
both $X$ and $X'$ is still bounded away from $0$:
\begin{claim}\label{claim:val-intersection-restated}
$\pE_{\mu \times \mu}[\val_I(X \wedge X')] \geq c^2$.
\end{claim}
\begin{align*}
\pE_{\mu \times \mu}[\val_I(X \wedge X')] &= \pE_{\mu \times \mu}[\E_{e \sim E(G)}[\one(e \text{ satisfied by } X)\one(e \text{ satisfied by } X')]] \\
&= \E_{e \sim E(G)}[\pPr_\mu[e \text{ satisfied by } X]^2] \\
&\geq \E_{e \sim E(G)}[\pPr_\mu[e \text{ satisfied by } X]]^2 ~~~~~~~~~~\text{(by Cauchy-Schwarz)}\\
&= \val_\mu(I)^2.
\end{align*}




%\mitali{write for general $J_{n,\ell,\alpha}$?}

\begin{lemma}\label{lem:sp-j-low-comp}
For all positive constants $c, \alpha, \tau \leq 1$, $\nu = \frac{c^4}{|\Sigma|^2\ell^{O(r)}}$, all integers $\ell \geq \Omega(r), n \geq \ell$, where $r = \Theta(\frac{\log c}{\log(1-\alpha)})$ the following holds: Let $I$ be an affine UG instance on $J_{n,\ell,\alpha}$ and $\mu$ be a pseudodistribution over assignments for $I$ with $\val_\mu(I) \geq c$ and degree at least $ \frac{\ell^{O(r)}\poly(|\Sigma|)}{\poly(c)\tau}+ \ell^{\poly(|\Sigma|\ell^r/c)}$. Then there exists a restriction $a \subseteq [n]$ of size $i \leq r$, a degree $\widetilde{O}(1/\nu)$ polynomial $P_a(X,X')$ in a fixed set of $|\Sigma|$ polynomials, subsets $A, B \subseteq V(J|_a)$ of size at most $\frac{\ell^{O(i)}|\Sigma|^2}{c^2\tau}$ and strings $y_A, y_B$ such that conditioning $\mu$ on the events $Y_A = y_A$ and $Y_B = y_B$ gives degree $\widetilde{\Omega}(1/\nu)$ pseudodistributions $\mu_1$ and $\mu_2$ such that:
\begin{enumerate}
\item $\cA_I  ~~\vdash_{\widetilde{O}(1/\nu)} ~~ P_a(X,X') \in [0,1]$.
\item $\pE_{\mu_1 \times \mu_2}[\Phi^a_{\beta,\nu}(X,X')P_a(X,X')] \geq \frac{c^{\Omega(r)}}{\exp(r^2)} \pE[P_a(X,X')]$.
\item $\pE_{\mu_1 \times \mu_2}[P_a(X,X')] \geq \Omega\left(\frac{c^2}{8r\ell^i|\Sigma|}\right)$.
\item $\E_{\substack{u_1,v_1 \sim S: u_1 \neq v_1\\ u_2,v_2 \sim S: u_2 \neq v_2}}[I(Y_{u_1,v_1};Y_{u_2,v_2})] \leq \tau.$
\item $\E_{\substack{u_1,v_1 \sim S: u_1 \neq v_1\\ u_2,v_2 \sim S: u_2 \neq v_2}}[I(Y'_{u_1,v_1};Y'_{u_2,v_2})] \leq \tau,$
\end{enumerate}
where $\beta = \left(\frac{\poly(c)}{\exp(r)}\right)^{r-i+1}$, $S = V(J|_a)$, $Y_{u,v} = (X_u,X_v,p_u,p_v)$, $Y'_{u,v} = (X'_u,X'_v,p'_u,p'_v)$ and the mutual information is taken with respect to the collection of local distributions $\cL(\mu_1 \times \mu_2,p_{\beta,\nu})$. 
\end{lemma}

\begin{proof}
Our strategy will be to use the edge-covering theorem. We will apply $\pE_{\mu \times \mu}$ on both sides of the inequality and use Claim~\ref{claim:val-intersection-restated} to say that $\pE[\val_I(X \wedge X')]$ is large. 

We will set parameters in the end of the proof such that: $r$ is large enough so that $4(1-\alpha)^{r+1} \leq c^2/6$, $\nu$ is small enough such that $|\Sigma|\ell^{O(r)}\sqrt{\nu} \leq c^2/6$ and $\eps_i$'s satisfy $2^{6r}\max_{i}\frac{\eps_{i-1}}{\eps_i^4} \leq c^2/6$ and further $\eps_i$'s satisfy the conditions of Theorem~\ref{thm:edge-covering-restated}. So we get that the error term in Theorem~\ref{thm:edge-covering-restated} is small, 
\begin{equation}\label{eq:error}
{\sf err} = 4(1-\alpha)^{r+1} + 2^{6r}\max_{i}\frac{\eps_{i-1}}{\eps_i^4}+ |\Sigma|\ell^{O(r)}\sqrt{\nu} \leq \frac{c^2}{2}.
\end{equation}

Applying $\pE_{\mu \times \mu}$ on both sides of Theorem~\ref{thm:edge-covering-restated} (assuming $\deg(\mu) \geq \ell^{\poly(\ell^r/\nu)}$) and 
using Claim~\ref{claim:val-intersection-restated} to lower bound pseudoexpectation of the LHS we get:
\[
c^2
\leq
\pE_{\mu \times \mu}[\val_I(X \wedge X')]
\leq \sum_{i \in [r]} \pE_{\mu \times \mu}[T_i(X,X')]+{\sf err}
\leq \sum_{i \in [r]} \pE_{\mu \times \mu}[T_i(X,X')] + \frac{c^2}{2},\]
% so that $\sum_{i \in [r]} \pE[T_i(X,X')]\geq \frac{3 c^2}{4}$.
% For an $i$-restriction $a \in {[n] \choose i}$ let $R_{s,a}$ denote the expression $P_{\eps_i,\nu}(\delta(G_s|_a)) (1-\prod_{j < i}\prod_{b \subset a: |b| = j}(1-P_{\eps_j,\nu}(\delta(G_s|_b))))$. Note that if the functions $P_{\beta,\nu}$ were exact indicators this denotes the event where $\mu(G_s|_a) \geq \eps_i$, but for \emph{all} smaller restrictions $b \subset a$, $\delta(G_s|_b) < \eps_{|b|}$. 
Therefore we get that there is an $i \in [r]$ for which $\pE_{\mu \times \mu}[T_i(X,X')] \geq c^2/2r$, and so
\begin{equation}\label{eq:low-comp1}
\pE_{\mu \times \mu}\E_{a \sim {[n] \choose i}}[\sum_s R_{s,a} \E_{u \sim J|_a}[G_s(u)\val^a_u(X \wedge X')]] \geq \frac{c^2}{2r\ell^i}.
\end{equation}

We aim to condition $\pE_{\mu \times \mu}$ on an appropriate polynomial (as in the proof of Lemma~\ref{lem:sp-j-restated}) so that the resulting pseudoexpectation operator satisfies: $\pE'[\E_{u \sim J|_a}[G_s(u)\val^a_u(X \wedge X')]]] \geq \Omega(c^2/r)$ for some specific $i$-restriction $a$. To do so we apply $\pE_{\mu \times \mu}$ to Claim~\ref{claim:dense-subcubes} and sum up over $s$, to get
% \begin{claim}\label{claim:dense-subcubes}
% If $\eps_0, \ldots, \eps_r$ satisfy that $\eps_{i-1} \leq \eps_i/(2^{i+1} i)$ then we get that:
% \[\vdash_{\tO(2^i/\nu)} \E_{a \sim {[n] \choose i}}[R_{s,a}] \leq \frac{4\delta(G_s)}{\eps_i^2\ell^i}.\]
% \end{claim}
% Summing up over $s \in \Sigma$ and applying $\pE_{\mu \times \mu}$ in the above claim we get:
\[\pE_{\mu \times \mu}\left[\E_{a \sim {[n] \choose i}}\sum_s R_{s,a}\right] \leq \frac{4}{\eps_i^2 \ell^i} + O(|\Sigma|\nu) \leq \frac{5}{\eps_i^2\ell_i},\]
where for the last inequality we have assumed $\nu \leq O(1/\ell^i|\Sigma|)$.

We split the RHS in equation~\ref{eq:low-comp1} to $c^2/4r\ell^i + c^2/4r\ell^i$, where in the first term we further use the above equation: $1/\ell^i \geq \frac{\eps_i^2}{5}\pE_{\mu \times \mu}[\E_{a \sim {[n] \choose i}}[\sum_s R_{s,a}]]$ to get that,
\begin{equation*}
\pE_{\mu \times \mu}\E_{a \sim {[n] \choose i}}\left[\sum_s R_{s,a} \E_{u \sim J|_a}[G_s(u)\val^a_u(X \wedge X')]\right] \geq \frac{c^2\eps_i^2}{20r}\pE_{\mu \times \mu}\left[\E_{a \sim {[n] \choose i}}[\sum_s R_{s,a}]\right] + \frac{c^2}{4r\ell^i},
\end{equation*}
and by re-arranging 
\begin{equation*}
\pE_{\mu \times \mu}\E_{a \sim {[n] \choose i}}\left[\sum_s R_{s,a} \left(\E_{u \sim J|_a}[G_s(u)\val^a_u(X \wedge X')] - \frac{c^2\eps_i^2}{20r}\right)\right] \geq \frac{c^2}{4r\ell^i}.
\end{equation*}
By averaging we get an $a \in {[n] \choose i}, s \in \Sigma$ such that:
\begin{equation}\label{eq:low-comp2}
\pE_{\mu \times \mu}\left[R_{s,a} \left(\E_{u \sim J|_a}[G_s(u)\val^a_u(X \wedge X')] - \frac{c^2\eps_i^2}{20r}\right)\right] \geq \frac{c^2}{4r\ell^i|\Sigma|}.
\end{equation}
The last result is the same as equation~\eqref{eq:sp-j-2}
in the proof of Lemma~\ref{lem:sp-j-restated}, 
which is the analogous statement for the completeness s clsoe to $1$ case. The rest of the proof is essentially the same as the rest of the proof of Lemma~\ref{lem:sp-j-restated}, and is outlined below. 

Applying Lemma~\ref{lem:ragh-tan-appln} we get pseudodistributions $\mu_1,\mu_2$ with low global correlation (as required in points 4,5 of the Lemma) with:
\begin{equation}\label{eq:low-comp3}
\pE_{\mu_1 \times \mu_2}\left[R_{s,a} \left(\E_{u \sim J|_a}[G_s(u)\val^a_u(X \wedge X')] - \frac{c^2\eps_i^2}{20r}\right)\right] \geq \frac{c^2}{8r\ell^i|\Sigma|}.
\end{equation}

Letting $P_a(X,X') = R_{s,a}(X,X')$ it is easy to check that:
\[\cA_I \vdash_{\tO(\exp(i)/\nu)} P_a(X,X') \in [0,1], \text{ and }~~~ \pE_{\mu_1 \times \mu_2}[P_a(X,X')] \geq \frac{c^2}{8r\ell^i|\Sigma|},\]
which proves points 1 and 3 of the Lemma statement. The only thing remaining to check is that conditioning $\mu_1 \times \mu_2$ on $R_{s,a}$ results in large shift-partition potential on $J|_a$. Recall that the shift-partition potential on $J|_a$ is: $\phi^a_{\beta,\nu}(X,X') = \sum_s \E_{u \sim J|_a}[G_s(u)p_{\beta,\nu}(\val^a_u(X))p_{\beta,\nu}(\val^a_u(X'))]^2$. From equation~\eqref{eq:low-comp3} we have that conditioning $\mu_1 \times \mu_2$ on $R_{s,a}$ gives:
\[\pE_{\mu_1 \times \mu_2 | R_{s,a}}[\E_{u \sim J|_a}[G_s(u)\val^a_u(X \wedge X')]] \geq  \frac{c^2\eps_i^2}{20r}.\]

We can prove that for $x \in [0,1] \vdash_{\tO(1/\nu)} x \leq  p_{\beta,\nu}(x) + \beta + O(\nu)$ (using the standard machinery in Section~\ref{sec:apx-ind}). Applying this with $\val^a_u(X \wedge X')$ and $\beta = \frac{c^2\eps_i^2}{40r}$, and further noting that
$p_{\beta,\nu}(\val_u^a(X \wedge X')) \leq p_{\beta,\nu}(\val_u^a(X))p_{\beta,\nu}(\val_u^a(X')) + O(\nu)$ we get:
\[\pE_{\mu_1 \times \mu_2 | R_{s,a}}[\E_{u \sim J|_a}[G_s(u)p_{\beta,\nu}(\val_u^a(X))p_{\beta,\nu}(\val_u^a(X'))]] \geq  \frac{c^2\eps_i^2}{40r} - O(\nu).\]

Applying Cauchy-Schwarz we get:
\[\pE_{\mu_1 \times \mu_2 | R_{s,a}}[\E_{u \sim J|_a}[G_s(u)p_{\beta,\nu}(\val_u^a(X))p_{\beta,\nu}(\val_u^a(X'))]^2] \geq  \Omega\left(\frac{c^4\eps_i^4}{r^2}\right).\]

By adding the terms corresponding to other $s \in \Sigma$ we immediately get:
\begin{equation}\label{eq:potential-c0}
\pE_{\mu_1 \times \mu_2 | R_{s,a}}[\phi_{\beta,\nu}^a(X,X')] \geq \Omega\left(\frac{c^4\eps_i^4}{r^2}\right).
\end{equation}

We finish the proof by setting $r,\eps_i,\nu$ according to all the requirements above.

\paragraph{Setting parameters:} To make $\sf{err}$ in \eqref{eq:error} small we set $r = \Theta(\frac{\log c}{\log(1-\alpha)})$, $\nu = \frac{c^4}{|\Sigma|^2\ell^{O(r)}}$, $\eps_{j-1} = \eps_j^5/2^{6r}$ with $\eps_r = \Theta(c^2\exp(-r))$. We need $\beta = c^2\eps_i^2/40r$ which equals $(\poly(c)/\exp(r))^{r-i+1}$. We can set $\deg(\mu) = \ell^{\poly(\ell^r/\nu)} = \ell^{\poly(|\Sigma|\ell^r/c)}$ so that it is valid to apply the edge-covering theorem. Additionally when we apply Lemma~\ref{lem:ragh-tan-appln} we need $\deg(\mu) \geq \Omega(\log |\Sigma|/p\tau\nu)$ where $p = c^2/4r\ell^i|\Sigma|$ so it suffices to take $\deg(\mu) = \frac{\ell^{O(r)}\poly(|\Sigma|)}{\poly(c)\tau}$. One can check that these parameters are good enough to make the rest of the statements in the proof work out. From \eqref{eq:potential-c0} we get that the conditioned shift-partition potential is at least $\Omega(c^2\eps_i^4/r^2) \geq \frac{c^{O(r)}}{\exp(r)}$ and this completes the proof.
\end{proof}


\subsection{Proof of Theorem~\ref{thm:main_johnson}}
The algorithm for arbitrary completeness is the same as that for completeness close to $1$ albeit with different parameter settings for $r, \delta,D,\gamma$ for the main algorithm and $t,\tau,\beta,\nu$ for SubRound. Given Lemma~\ref{lem:sp-j-low-comp} the rest of the analysis of the algorithm remains exactly the same. We set these parameters below but omit the description of the algorithm. The guarantees of this algorithm are potentially worse (running time $\approx n^{\ell^{\poly(\ell)}}$ and approximation factor $\approx c^{O(\log 1/c)}$) than the one for completeness close to $1$, but the important thing to note is that the running time is polynomial in $n$ (as long as $\ell, |\Sigma|$ are $O(1)$) and the approximation factor only depends on $c$ and not on $\ell,n$ or $|\Sigma|$.

\begin{theorem}\label{thm:main-johnson-arbit-comp}
For all positive constants $c \in (0,1]$, $\alpha \in \Q$, all integers $\ell \geq \Omega(r)$ with $\alpha \ell \in \N$ and $r =  \Theta(\frac{\log c}{\log(1-\alpha)})$ and $n$ large enough, Algorithm~\ref{alg:j} (with modified parameter settings given below) has the following guarantee: If $I$ is an instance of affine Unique Games on the $(n,\ell,\alpha)$-Johnson graph $J$ with alphabet $\Sigma$ and $\val(I) = c$, then in time $|V(J)|^{D}$ Algorithm~\ref{alg:j} returns an $\Omega\left(\frac{c^{\Omega(r)}}{\exp(r^2)}\right)$-satisfying assignment for $I$, with $D = \frac{\ell^{\poly(|\Sigma|\ell^r/c)}}{c^{O(r)}}$.
\end{theorem}



\begin{proof}
Let $r = \Theta(\frac{\log c}{\log(1-\alpha)})$. Given the $i \leq r$-restricted subcube $a \subseteq [n], P_a(X,X')$ and $\mu_1 \times \mu_2$ from Lemma~\ref{lem:sp-j-low-comp} we have that $\Phi_{\beta,\nu}^a(\mu_1 \times \mu_2| P_a) \geq \frac{c^{\Omega(r)}}{\exp(r^2)}$, therefore we apply Lemma~\ref{lem:correlation} with $S = J|_a$, the polynomial  $E(X,X') = P_a(X,X')$ and the parameters $\beta = \left(\frac{\poly(c)}{\exp(r)}\right)^{r-i+1}, \nu = \frac{c^4}{|\Sigma|^2\ell^{O(r)}}, \delta = \frac{c^{O(r)}}{\exp(r^2)}, p = \Omega\left(\frac{c^2}{8r\ell^i|\Sigma|}\right)$ and $\tau = \frac{c^{O(r)}}{\exp(r^2)\ell^{O(i)}|\Sigma|}$. The parameter $\tau$ has been chosen so that $O(\frac{\sqrt{\tau}+1/|S|}{p\delta^2}) = \frac{c^{O(r)}}{\exp(r^2)}$ and therefore we can apply the rounding lemma (Lemma~\ref{lem:round-j}) with $\zeta = \frac{c^{O(r)}}{\exp(r^2)}$ and $\gamma = \frac{c^{\Omega(r)}}{\exp(r^2)}$ and the same settings of $\beta, \nu, \delta$. This shows that there exists a subcube $a \in {[n] \choose i}$ with SubRound value that is at least $\frac{c^{\Omega(r)}}{\exp(r^2)}$ if $\deg(\mu) \geq  \frac{\ell^{O(r)}\poly(|\Sigma|)}{\poly(c)\tau}+ \ell^{\poly(|\Sigma|\ell^r/c)}$. Hence it suffices to have degree of $\mu$ equal to $D = \frac{\ell^{\poly(|\Sigma|\ell^r/c)}}{c^{O(r)}}$.

By using SubRound as a subroutine, we finish the proof of this theorem by applying the iteration Lemma~\ref{lem:subroutine} with completeness $c$, $\gamma = c^2$, $\delta = \frac{c^{\Omega(r)}}{\exp(r^2)}$ and $r = \Theta(\frac{\log c}{\log(1-\alpha)})$ so that $(1-\alpha)^r = \poly(c)$. We can check that SubRound (with the parameter settings above) satisfies the hypotheses of the lemma: it finds a $\le r$-restricted subcube and an assignment to it with value $\Omega(\delta)$ in time $|V(J)|^{D}$ (where $D$ is the degree bound above). This gives us that Algorithm~\ref{alg:j} outputs an assignment of value at least $\Omega\left(\frac{c^{\Omega(r)}}{\exp(r^2)}\right)$ in time $|V(G)|^{O(D)}$.
\end{proof}






















